\section{Conclusiones}
    Realizar este proyecto, y específicamente al ver los resultados de las 
    transformaciones, me permitió interiorizar porque el preprocesamiento de 
    las imágenes es una etapa tan más relevante al momento de hacer ciencia 
    de datos o entrenar modelos de IA, ya que una imagen se puede transformar 
    en una amplia variedad de versiones en las que se pueden resaltar de mejor 
    manera ciertos elementos o generalizarlos, según el fin del preprocesamiento. 
    Por esto pienso que entender cómo una imagen es modificada por una transformación 
    permitirá hacer una selección más cuidadosa sobre cuáles aplicar de forma 
    que se alinee con el objetivo del proyecto, por ejemplo, si una imagen está 
    conformada por elementos que solo pueden estar en una única orientación para 
    interpretarlos correctamente (dígase letras, números, rasgos faciales, estructuras 
    químicas), las transformaciones que podemos aplicar se encuentran limitadas 
    debido a que cualquier mínima alteración podría evocar a conclusiones erróneas.  
    Otro detalle que me percate es que cualquier transformación genera una pérdida de 
    información sobre lo que había originalmente en la imagen, haciendo que las 
    transformaciones no sean operaciones ideales (o totalmente reversibles), al menos 
    que se cuente con una caja de visión lo suficientemente grande (es decir, el tamaño 
    de la gráfica sea más grande lo que realmente abarca la imagen). \\

    Este miniproyecto fue más de implementación que otra actividad, pero ayudó a 
    entender cómo algunos softwares de edición de imagen trabajan a bajo nivel, 
    así como también expandir estas ideas del preprocesamiento y tratamiento 
    de imágenes previo a una análisis o estudio donde se tenga mayormente 
    datos visuales. Por ello, este proyecto abre las puertas a técnicas 
    auxiliares y complementarios totalmente útiles en el labor de la ciencia 
    de datos